\documentclass[11pt]{article}
\usepackage[T1]{fontenc}
\usepackage[utf8]{inputenc}
\usepackage[spanish,es-noshorthands]{babel}
\usepackage{lmodern}
\usepackage[margin=2.4cm]{geometry}
\usepackage{amsmath,amssymb}
\usepackage{booktabs}
\usepackage{xcolor}
\usepackage{enumitem}
\usepackage{parskip}
\usepackage[colorlinks=true,linkcolor=black,urlcolor=blue]{hyperref}

\definecolor{acmuerta}{HTML}{9AA0A6}
\definecolor{aclatente}{HTML}{C9A227}
\definecolor{acactiva}{HTML}{2E7D32}

\newcommand{\estado}{S}
\newcommand{\campo}{E}
\newcommand{\MUERTA}{\textcolor{acmuerta}{\textsf{MUERTA}}}
\newcommand{\LATENTE}{\textcolor{aclatente}{\textsf{LATENTE}}}
\newcommand{\ACTIVA}{\textcolor{acactiva}{\textsf{ACTIVA}}}

\title{\textbf{AstroBioSim — Modelo matemático del autómata celular}\\[2pt]
\large Decisiones de diseño para la exposición y defensa final}
\author{Área de Matemática}
\date{Julio 2026}

\begin{document}
\maketitle

\begin{abstract}
\noindent AstroBioSim simula la supervivencia poblacional de microorganismos en
subsuelos de Tierra, Marte y Encelado mediante un \textbf{autómata celular (AC)}
estocástico y espacialmente explícito sobre una grilla 2D. Este documento define
formalmente el AC: su espacio de estados de \emph{tres} niveles, la vecindad de
Moore, la actualización síncrona, la cinética de crecimiento continua y las
reglas de transición intercambiables. Es la pieza matemática central del proyecto
y el eje de la pregunta de investigación.
\end{abstract}

\section{Definición formal del autómata}

Un autómata celular es una tupla $\mathcal{A}=(\mathcal{L},\ \mathcal{Q},\
\mathcal{N},\ \phi)$ donde $\mathcal{L}$ es el retículo, $\mathcal{Q}$ el conjunto
de estados, $\mathcal{N}$ la vecindad y $\phi$ la función de transición local.

\paragraph{Retículo.} $\mathcal{L}=\{(i,j): 0\le i<M,\ 0\le j<N\}$, una grilla
$M\times N$. Cada celda es un \emph{sitio} físico del subsuelo, no un individuo:
las clases del programa son las \emph{especies} y los \emph{entornos}; las celdas
son entradas de arreglos \textsf{NumPy}.

\paragraph{Espacio de estados de tres niveles.} A diferencia del Juego de la Vida
clásico (vivo/muerto), aquí
\[
\mathcal{Q}=\{\,\MUERTA=0,\ \ \LATENTE=1,\ \ \ACTIVA=2\,\}.
\]
La distinción es la decisión matemática de fondo del proyecto (ADR-0012). Responde
a un hecho biológico: los extremófilos sobreviven la desecación en un estado
\emph{latente} (vivos, sin reproducirse) que un modelo binario no puede
representar. Formalmente:
\begin{itemize}[nosep]
  \item \ACTIVA: la célula se reproduce ($\mu>0$).
  \item \LATENTE: viva pero sin reproducirse ($\mu=0$); \emph{ocupa} el sitio.
  \item \MUERTA: sitio sin biomasa viable. Es un \textbf{estado absorbente}:
  la muerte es irreversible. Repoblar un sitio \MUERTA{} es \emph{colonización}
  por un vecino \ACTIVA, no resurrección de la misma célula.
\end{itemize}

\paragraph{Vecindad de Moore.} $\mathcal{N}(i,j)$ son los $8$ sitios adyacentes
(ortogonales y diagonales). Definimos los conteos de vecinos
\begin{align}
n_A(i,j) &= \#\{(k,l)\in\mathcal{N}(i,j): \estado_{k,l}=\ACTIVA\},\\
n_O(i,j) &= \#\{(k,l)\in\mathcal{N}(i,j): \estado_{k,l}\in\{\ACTIVA,\LATENTE\}\}.
\end{align}
La \textbf{condición de borde} por defecto es \emph{frontera muerta}: los sitios
fuera de la grilla cuentan como \MUERTA. Es coherente con que la grilla es una
\emph{ventana} del subsuelo, no un mundo cerrado. Se ofrece también la variante
\emph{toroidal} (wrap-around) para estudios sin efecto de borde.

\paragraph{Actualización síncrona (doble buffer).} El estado en $t{+}1$ se calcula
\emph{íntegro} a partir del estado en $t$:
\[
\estado^{\,t+1}_{i,j}=\phi\!\left(\estado^{\,t}_{i,j},\ \{\estado^{\,t}_{k,l}\}_{\mathcal{N}},\ \campo_{i,j}\right).
\]
Nunca se actualiza la grilla \emph{in situ}: se lee de $\estado^{\,t}$ y se escribe
en un buffer nuevo. Esto garantiza que el resultado es independiente del orden de
barrido, condición necesaria para que la vectorización \textsf{NumPy} sea
\emph{equivalente} a la definición celda-a-celda. Prohibimos los bucles Python
sobre celdas: todo se expresa con máscaras booleanas y desplazamientos de arreglos.

\section{Cinética de crecimiento continua (ADR-0013)}

El corazón cuantitativo es reemplazar la máscara binaria ``habitable / no
habitable'' por una \textbf{tasa de crecimiento} $\mu$ que depende de cuán lejos
del óptimo está el ambiente. Adoptamos el \emph{gamma concept} de la microbiología
predictiva: factores multiplicativos independientes, cada uno en $[0,1]$,
\begin{equation}
\mu(T,a_w,\mathrm{UV})=\mu_{\mathrm{opt}}\cdot\gamma_T(T)\cdot\gamma_{a_w}(a_w)\cdot\gamma_{\mathrm{UV}}(\mathrm{UV}).
\end{equation}

\paragraph{Factor térmico (modelo cardinal con inflexión, CTMI).} Con los tres
puntos cardinales $T_{\min}<T_{\mathrm{opt}}<T_{\max}$ y sin parámetros libres
(Rosso, Lobry \& Flandrois, 1993):
\begin{equation}
\gamma_T(T)=\frac{(T-T_{\max})(T-T_{\min})^2}
{(T_{\mathrm{opt}}-T_{\min})\big[(T_{\mathrm{opt}}-T_{\min})(T-T_{\mathrm{opt}})-(T_{\mathrm{opt}}-T_{\max})(T_{\mathrm{opt}}+T_{\min}-2T)\big]}
\end{equation}
para $T\in(T_{\min},T_{\max})$, y $0$ fuera. Se verifica $\gamma_T(T_{\mathrm{opt}})=1$
y $\gamma_T(T_{\min})=\gamma_T(T_{\max})=0$.

\paragraph{Factores de agua y de UV.}
\begin{equation}
\gamma_{a_w}(a_w)=\left[\frac{a_w-a_{w,\min}}{1-a_{w,\min}}\right]_0^1,\qquad
\gamma_{\mathrm{UV}}(\mathrm{UV})=\left[1-\frac{\mathrm{UV}}{\mathrm{UV}_{\max}}\right]_0^1,
\end{equation}
donde $[\,\cdot\,]_0^1$ denota recorte a $[0,1]$. Ambos se anulan al cruzar el
umbral de crecimiento respectivo.

\paragraph{Probabilidad de reproducción por tick.} La tasa continua se convierte en
la probabilidad de que una celda \ACTIVA{} coloque una hija en un sitio vecino
vacío durante un paso de duración $\Delta t$:
\begin{equation}
p_{\mathrm{rep}}=\big[\mu\cdot\Delta t\big]_0^1,\qquad \Delta t = 1\ \text{h}.
\end{equation}
Con $\mu_{\mathrm{opt}}$ en $\mathrm{h}^{-1}$, el producto es adimensional. Elegimos
$\Delta t=1$~h porque produce un gradiente informativo entre especies en el óptimo
($p_{\mathrm{rep}}\approx 0{,}88$ para \emph{E.\ coli}, $0{,}23$ para
\emph{D.\ radiodurans}, $0{,}07$ para \emph{M.\ burtonii}), en lugar de saturar a
$1$ como haría un paso diario.

\section{Reglas de transición intercambiables (ADR-0016)}

La función de transición $\phi$ no está cableada: es una \textbf{estrategia
intercambiable}. El motor calcula una vez por tick las máscaras ambientales
$\mathrm{cre}=\gamma\text{-crecimiento}$ y $\mathrm{sup}=$ supervivencia, la
probabilidad $p_{\mathrm{rep}}$ y los conteos $n_A,n_O$, y se los entrega a la
regla. Cada regla se autodescribe en notación formal para el panel de la interfaz.
Sea $U\sim\mathcal{U}(0,1)$ una variable del generador sembrado (reproducibilidad).

\paragraph{Regla logística (proceso de contacto) — por defecto.}
\begin{equation}
\estado^{\,t+1}_{i,j}=
\begin{cases}
\MUERTA & \neg\,\mathrm{sup}(\campo_{i,j})\\[2pt]
\ACTIVA & \text{ocupada}\ \wedge\ \mathrm{cre}(\campo_{i,j})\\[2pt]
\LATENTE & \text{ocupada}\ \wedge\ \mathrm{sup}\ \wedge\ \neg\,\mathrm{cre}\\[2pt]
\ACTIVA & \text{vacía}\ \wedge\ \mathrm{cre}\ \wedge\ U<p_{\mathrm{rep}}\,\dfrac{n_A}{8}\\[6pt]
\MUERTA & \text{en otro caso}
\end{cases}
\end{equation}
Sin muerte por soledad; la única muerte es ambiental. Es la más defendible
biológicamente: un frente de colonización avanza sobre sitios viables a una
velocidad regida por $\mu$.

\paragraph{Regla de Conway (B3/S23), filtrada por el ambiente.}
\begin{equation}
\estado^{\,t+1}_{i,j}=
\begin{cases}
\MUERTA & \neg\,\mathrm{sup}\\
\{\ACTIVA,\LATENTE\} & \text{ocupada}\ \wedge\ n_O\in\{2,3\}\\
\ACTIVA & \text{vacía}\ \wedge\ \mathrm{cre}\ \wedge\ n_O=3\ \wedge\ U<p_{\mathrm{rep}}\\
\MUERTA & \text{en otro caso}
\end{cases}
\end{equation}
Se conserva por su valor visual y como validación (bajo ambiente óptimo reproduce
el Juego de la Vida: un \emph{blinker} oscila con periodo $2$, lo que testeamos).
Los umbrales $2$–$3/3$ son del autómata, no de la microbiología.

\paragraph{Regla híbrida.} Contacto para nacer más una muerte por sobrepoblación
($n_O>k$) que evita la saturación total, sin muerte por soledad. El editor de
reglas por bloques de la interfaz (Hito 3) produce nuevas reglas sobre esta misma
firma, sin tocar el motor.

\section{Reproducibilidad y ensambles Montecarlo}

Toda la aleatoriedad (reproducción, eventos estocásticos) usa un único
\textsf{numpy.random.Generator} con semilla explícita, inyectado como parámetro:
nunca el estado global. En consecuencia, \emph{misma semilla $\Rightarrow$ misma
corrida}. Como el sistema es estocástico, una sola corrida no es la respuesta: se
promedian $R$ repeticiones con semillas $s_0,s_0{+}1,\dots$ y se reporta la
media $\pm$ desviación de las fracciones de sitios \ACTIVA/\LATENTE/\MUERTA a lo
largo del tiempo. Validamos que $R$ es suficiente cuando la media se estabiliza al
crecer $R$.

\section{La pregunta de investigación (ADR-0015)}

El aporte del proyecto es cuantitativo. En el regolito marciano en bulk
($a_w\approx 0{,}19$, muy por debajo del límite de división $\approx 0{,}605$)
nadie crece: el resultado correcto es la extinción del crecimiento activo. La vida
persiste solo en \emph{microrefugios} húmedos transitorios (salmueras). Formulamos:

\begin{quote}\itshape
¿Con qué frecuencia y magnitud mínimas deben aparecer microrefugios húmedos
transitorios para que una población persista en un subsuelo planetario, en vez de
extinguirse?
\end{quote}

Operativamente es un \textbf{barrido de parámetros} $\text{frecuencia}\times
\text{magnitud}$ del evento de salmuera sobre ensambles Montecarlo, con la
persistencia (fracción no-\MUERTA{} al final) como variable de respuesta. La salida
es un \emph{umbral crítico}: por debajo la población se extingue, por encima
persiste. La combinación espacialidad $+$ estocasticidad $+$ latencia es un nicho
poco cubierto por la literatura de habitabilidad, que suele ser puntual (0-D).

\section*{Síntesis de decisiones (defensa)}
\begin{center}
\small
\begin{tabular}{@{}ll@{}}
\toprule
\textbf{Decisión} & \textbf{Justificación} \\
\midrule
Tres estados (no binario) & representar la latencia sin violar la muerte irreversible (ADR-0012)\\
Cinética CTMI continua & $\mu$ deja de ser decorativa; ``crece a tal tasa'' en vez de ``vive/muere'' (ADR-0013)\\
Reglas intercambiables & menú de reglas + editor por bloques + panel de notación (ADR-0016)\\
Frontera muerta & la grilla es una ventana del subsuelo, no un toro\\
$\Delta t = 1$~h & gradiente informativo de $p_{\mathrm{rep}}$ entre especies\\
Semilla inyectada & reproducibilidad exacta; base de los ensambles Montecarlo\\
\bottomrule
\end{tabular}
\end{center}

\end{document}
